% GENERATED FILE. Edit canonical lessons, then run npm run generate:books.
% canonical-source-hash: fnv1a64:e7483b0025c2a853
% canonical-lessons: MR-C400-vahan, MR-C400-ducaki, MR-C400-riksa, MR-C400-trak, MR-C400-kombne

\chapter{Two Wheels and Three}
\label{ch:mr-two-wheels-and-three}
\hlchaptermodality{400}

\begin{chapteropening}
\textbf{By the end of this chapter:} \emph{I can talk about vehicles, scooters, autorickshaws and trucks.}

Complete the last lesson of chapter 400: I can talk about vehicles, scooters, autorickshaws and trucks.
\end{chapteropening}

\section[\texorpdfstring{\mr{वाहन} (vāhan)}{वाहन (vāhan)}]{\mr{वाहन} (vāhan) — a vehicle}

\label{lesson:MR-C400-vahan}



\begin{quote}
Before the new one: say the Marathi for a fellow traveller, then the Marathi for a wayfarer.
\end{quote}

\subsection*{You'll want to know: \mr{वाहन}}
\textbf{\mr{वाहन}} — \emph{vāhan} — \textquotedblleft{}a vehicle\textquotedblright{}.

\begin{quote}
\mr{या} \mr{रस्त्यावर} \mr{मोठ्या} \mr{वाहनांना} \mr{बंदी} \mr{आहे}.
\end{quote}

\begin{quote}
\emph{yā rastyāvar moṭhyā vāhanānnā bandī āhe.}
\end{quote}

\begin{quote}
Large vehicles are banned on this road.
\end{quote}

Neuter.

A formal, general word; the plural before an ending is \mr{वाहनां}-: \mr{वाहनांना}.

\begin{cousinweb}
\textbf{\mr{वाहन}} (\emph{vāhan}) is a tatsama, taken from Sanskrit \emph{vāhana}, \textquotedblleft{}conveying; a vehicle, a mount\textquotedblright{} (Monier-Williams).
\end{cousinweb}

\begin{grammarlens}[title={What you've built}]
One more word for getting around.
\end{grammarlens}

\subsection*{Your turn}
\begin{itemize}
\raggedright
  \item \emph{Say it:} \emph{vāhan}
  \item \emph{Say it:} \emph{vāhan}, once more
  \item \emph{Say it:} \emph{yā rastyāvar moṭhyā vāhanānnā bandī āhe.} — Large vehicles are banned on this road.
  \item \emph{Recall:} say the Marathi for a fellow traveller, then the Marathi for a wayfarer, then say \emph{vāhan} again
\end{itemize}

\begin{tcolorbox}[breakable,colback=teal!4,colframe=teal!35!black,title={Before you move on}]
What does \mr{वाहन} mean? (\textquotedblleft{}A vehicle\textquotedblright{}.) Say it once more.
\end{tcolorbox}
\section[\texorpdfstring{\mr{दुचाकी} (ducākī)}{दुचाकी (ducākī)}]{\mr{दुचाकी} (ducākī) — a two-wheeler}

\label{lesson:MR-C400-ducaki}



\begin{quote}
Before the new one: say the Marathi for a wayfarer, then the Marathi for a vehicle.
\end{quote}

\subsection*{You'll want to know: \mr{दुचाकी}}
\textbf{\mr{दुचाकी}} — \emph{ducākī} — \textquotedblleft{}a two-wheeler\textquotedblright{}.

\begin{quote}
\mr{मीरा} \mr{दुचाकीवरून} \mr{कार्यालयात} \mr{जाते}.
\end{quote}

\begin{quote}
\emph{mīrā ducākīvarūn kāryālayāt jāte.}
\end{quote}

\begin{quote}
Meera rides to the office on her two-wheeler.
\end{quote}

Feminine.

Any scooter or motorbike; a three-wheeler is a \mr{तीनचाकी}.

\begin{cousinweb}
\textbf{\mr{दुचाकी}} (\emph{ducākī}) is a Marathi compound of \emph{du-}, \textquotedblleft{}two\textquotedblright{}, and \textbf{\mr{चाक}}, \textquotedblleft{}a wheel\textquotedblright{}, inherited from Sanskrit \emph{cakra} (CDIAL 4538), with the adjective ending -\mr{ई}.
\end{cousinweb}

\begin{grammarlens}[title={What you've built}]
One more word for getting around.
\end{grammarlens}

\subsection*{Your turn}
\begin{itemize}
\raggedright
  \item \emph{Say it:} \emph{ducākī}
  \item \emph{Say it:} \emph{ducākī}, once more
  \item \emph{Say it:} \emph{mīrā ducākīvarūn kāryālayāt jāte.} — Meera rides to the office on her two-wheeler.
  \item \emph{Recall:} say the Marathi for a wayfarer, then the Marathi for a vehicle, then say \emph{ducākī} again
\end{itemize}

\begin{tcolorbox}[breakable,colback=teal!4,colframe=teal!35!black,title={Before you move on}]
What does \mr{दुचाकी} mean? (\textquotedblleft{}A two-wheeler\textquotedblright{}.) Say it once more.
\end{tcolorbox}
\section[\texorpdfstring{\mr{रिक्षा} (rikṣā)}{रिक्षा (rikṣā)}]{\mr{रिक्षा} (rikṣā) — an autorickshaw}

\label{lesson:MR-C400-riksa}



\begin{quote}
Before the new one: say the Marathi for a vehicle, then the Marathi for a two-wheeler.
\end{quote}

\subsection*{You'll want to know: \mr{रिक्षा}}
\textbf{\mr{रिक्षा}} — \emph{rikṣā} — \textquotedblleft{}an autorickshaw\textquotedblright{}.

\begin{quote}
\mr{स्टेशनवरून} \mr{आम्ही} \mr{रिक्षाने} \mr{घरी} \mr{आलो}.
\end{quote}

\begin{quote}
\emph{sṭeśanvarūn āmhī rikṣāne gharī ālo.}
\end{quote}

\begin{quote}
We came home from the station by autorickshaw.
\end{quote}

Feminine.

The vehicle takes -\mr{ने} for \textquotedblleft{}by\textquotedblright{}: \mr{रिक्षाने}, \mr{बसने}, by rickshaw, by bus.

\begin{cousinweb}
\textbf{\mr{रिक्षा}} (\emph{rikṣā}) is a loan from English \emph{rickshaw}, which English took from Japanese; in Marathi it names the three-wheeled autorickshaw.
\end{cousinweb}

\begin{grammarlens}[title={What you've built}]
One more word for getting around.
\end{grammarlens}

\subsection*{Your turn}
\begin{itemize}
\raggedright
  \item \emph{Say it:} \emph{rikṣā}
  \item \emph{Say it:} \emph{rikṣā}, once more
  \item \emph{Say it:} \emph{sṭeśanvarūn āmhī rikṣāne gharī ālo.} — We came home from the station by autorickshaw.
  \item \emph{Recall:} say the Marathi for a vehicle, then the Marathi for a two-wheeler, then say \emph{rikṣā} again
\end{itemize}

\begin{tcolorbox}[breakable,colback=teal!4,colframe=teal!35!black,title={Before you move on}]
What does \mr{रिक्षा} mean? (\textquotedblleft{}An autorickshaw\textquotedblright{}.) Say it once more.
\end{tcolorbox}
\section[\texorpdfstring{\mr{ट्रक} (ṭrak)}{ट्रक (ṭrak)}]{\mr{ट्रक} (ṭrak) — a truck; a lorry}

\label{lesson:MR-C400-trak}



\begin{quote}
Before the new one: say the Marathi for a two-wheeler, then the Marathi for an autorickshaw.
\end{quote}

\subsection*{You'll want to know: \mr{ट्रक}}
\textbf{\mr{ट्रक}} — \emph{ṭrak} — \textquotedblleft{}a truck; a lorry\textquotedblright{}.

\begin{quote}
\mr{रात्री} \mr{महामार्गावर} \mr{खूप} \mr{ट्रक} \mr{असतात}.
\end{quote}

\begin{quote}
\emph{rātrī mahāmārgāvar khūp ṭrak astāt.}
\end{quote}

\begin{quote}
There are lots of trucks on the highway at night.
\end{quote}

Masculine.

The plural is the same word: \mr{एक} \mr{ट्रक}, \mr{दहा} \mr{ट्रक}.

\begin{cousinweb}
\textbf{\mr{ट्रक}} (\emph{ṭrak}) is a loan from English \emph{truck}.
\end{cousinweb}

\begin{grammarlens}[title={What you've built}]
One more word for getting around.
\end{grammarlens}

\subsection*{Your turn}
\begin{itemize}
\raggedright
  \item \emph{Say it:} \emph{ṭrak}
  \item \emph{Say it:} \emph{ṭrak}, once more
  \item \emph{Say it:} \emph{rātrī mahāmārgāvar khūp ṭrak astāt.} — There are lots of trucks on the highway at night.
  \item \emph{Recall:} say the Marathi for a two-wheeler, then the Marathi for an autorickshaw, then say \emph{ṭrak} again
\end{itemize}

\begin{tcolorbox}[breakable,colback=teal!4,colframe=teal!35!black,title={Before you move on}]
What does \mr{ट्रक} mean? (\textquotedblleft{}A truck; a lorry\textquotedblright{}.) Say it once more.
\end{tcolorbox}
\section[\texorpdfstring{\mr{कोंबणे} (kombṇe)}{कोंबणे (kombṇe)}]{\mr{कोंबणे} (kombṇe) — to stuff in; to cram}

\label{lesson:MR-C400-kombne}



\begin{quote}
Before the new one: say the Marathi for an autorickshaw, then the Marathi for a truck.
\end{quote}

\subsection*{You'll want to know: \mr{कोंबणे}}
\textbf{\mr{कोंबणे}} — \emph{kombṇe} — \textquotedblleft{}to stuff in; to cram\textquotedblright{}.

\begin{quote}
\mr{त्याने} \mr{पिशवीत} \mr{सगळे} \mr{कपडे} \mr{कोंबले}.
\end{quote}

\begin{quote}
\emph{tyāne piśvīt sagḷe kapḍe komble.}
\end{quote}

\begin{quote}
He stuffed all the clothes into the bag.
\end{quote}

Past, with -\mr{ने} (he …): \mr{त्याने} \mr{कोंबलं}, \emph{tyāne komblaṁ}.

The container takes -\mr{त}: \mr{पिशवीत} \mr{कोंबणे}; people crammed into a rickshaw are \mr{कोंबलेले}.

\begin{cousinweb}
\textbf{\mr{कोंबणे}} (\emph{kombṇe}) is inherited: CDIAL lists Marathi \emph{kõbṇẽ}, \textquotedblleft{}to thump, thrust, force in, cram in\textquotedblright{}, under \emph{*kōbb}, \textquotedblleft{}strike\textquotedblright{} (CDIAL 3522).
\end{cousinweb}

\begin{grammarlens}[title={What you've built}]
One more word for getting around.
\end{grammarlens}

\subsection*{Your turn}
\begin{itemize}
\raggedright
  \item \emph{Say it:} \emph{kombṇe}
  \item \emph{Say it:} \emph{kombṇe}, once more
  \item \emph{Say it:} \emph{tyāne piśvīt sagḷe kapḍe komble.} — He stuffed all the clothes into the bag.
  \item \emph{Recall:} say the Marathi for an autorickshaw, then the Marathi for a truck, then say \emph{kombṇe} again
\end{itemize}

\begin{tcolorbox}[breakable,colback=teal!4,colframe=teal!35!black,title={Before you move on}]
What does \mr{कोंबणे} mean? (\textquotedblleft{}To stuff in; to cram\textquotedblright{}.) Say it once more.
\end{tcolorbox}
